\addchap{Einleitung}
\label{ch:index}

Luanti verwendet Lua-Skripte, um Modding zu ermöglichen. Dieses Buch
soll dir, beginnend mit den Grundlagen, beibringen, wie du deine eigenen
Mods erstellst. Jedes Kapitel konzentriert sich auf einen bestimmten
Teil der API und bringt dich schnell dahin, eigene Mods zu schreiben.

Das Buch ist auch online lesbar\footnote{\url{https://rubenwardy.com/minetest_modding_book/de/}}
und als HTML herunterladbar\footnote{\url{https://gitlab.com/rubenwardy/minetest_modding_book/-/releases}}.

\section*{Feedback und Beiträge}\label{feedback-und-beitruxe4ge}
\addcontentsline{toc}{section}{Feedback und Beiträge}

Hast du einen Fehler bemerkt oder möchtest du Feedback geben? Dann teile
es bitte mit:

\begin{itemize}
\tightlist
\item
  Erstelle ein Issue auf GitLab\footnote{\url{https://gitlab.com/rubenwardy/minetest_modding_book/-/issues}}.
\item
  Schreib im Forenthema\footnote{\url{https://forum.luanti.org/viewtopic.php?f=14&t=10729}}
  oder im deutschen Forenthema\footnote{\url{https://forum.luanti.org/viewtopic.php?t=26723}}.
\item
  Nimm Kontakt mit dem Autor auf\footnote{\url{https://rubenwardy.com/contact/}}.
\item
  Lust mitzumachen? Lies die README\footnote{\url{https://gitlab.com/rubenwardy/minetest_modding_book/-/blob/master/README.md}}.
\end{itemize}
